\chapter{Background and Related Work}

This chapter reviews the theoretical foundations and related work necessary for building a scene-aware dialogue system. It progresses from the high-level requirements of human-robot interaction down to the specific visual perception models, structured scene representations, and efficient multimodal language generation techniques utilized in this thesis.

\section{Social Robotics and Grounded Human-Robot Interaction (HRI)}

Social robots operate in shared human environments where users expect them to exhibit awareness of their physical and social surroundings. Recent integrations of Large Language Models (LLMs) into human-robot interaction (HRI) have significantly improved conversational fluency. However, as Kim et al.~\cite{kim_understanding_2024} and Atuhurra et al.~\cite{atuhurra_leveraging_2024} note, this fluency often raises user expectations: when a robot speaks naturally, users implicitly assume it also understands the physical context and social nuances of the conversation.

To prevent a mismatch between linguistic competence and actual situational awareness, the robot's dialogue must be grounded. In robotics, grounding denotes the explicit connection between linguistic expressions (e.g., ``the person on my left'') and the physical entities to which they refer. Unlike a purely text-based assistant, an embodied robot must continuously resolve these references against its changing perceptual state.

The Pepper robot is a prominent platform for HRI research due to its humanoid form, native tablet interface~\cite{softbank_peppers_tablet_2026}, and built-in speech recognition~\cite{softbank_alspeechrecognition_2026}. However, its onboard computational hardware is severely limited and insufficient for modern deep learning inference. Consequently, achieving near real-time scene understanding on Pepper requires a distributed architecture, where heavy visual and linguistic processing is offloaded to an external server while the robot handles interaction and sensory acquisition~\cite{reyes_near_2018, trinquet_future_2025, brad_retrieval-augmented_2025}.

\section{Visual Perception for Robotics}

The foundation of scene awareness is object detection. Traditionally, robotics has relied on Convolutional Neural Networks (CNNs), such as the YOLO family, which utilize anchor boxes and dense grid predictions. While fast, these models require heuristic post-processing like Non-Maximum Suppression (NMS). Recently, the architectural shift toward Detection Transformers (DETRs)~\cite{zhao_detrs_2024} has reframed detection as a direct set-prediction problem. By leveraging self-attention mechanisms~\cite{vaswani2023attentionneed} and bipartite matching during training, models like RT-DETR and RF-DETR eliminate NMS, providing state-of-the-art accuracy while maintaining the real-time latency required for interactive robotics.

Standard detectors operate on a closed vocabulary, mapping visual features to a fixed set of classes via a classifier head. This restricts natural dialogue, as the robot cannot recognize novel objects. Open-vocabulary models, such as OWL-ViT~\cite{minderer_simple_2022} and Grounding DINO~\cite{liu_grounding_2024}, break this barrier by projecting image patches and text prompts into a shared multimodal embedding space (often aligned via CLIP~\cite{pmlr-v139-radford21a}). This allows the detection of arbitrary objects based on text similarity, a crucial capability for open-ended scene-aware dialogue.

Detecting an object in a single frame is insufficient for conversation; the system must preserve object permanence. Multi-Object Tracking (MOT) algorithms associate detections across frames using geometric overlap (IoU) and appearance embeddings~\cite{bewley_simple_2016, wojke_simple_2017, zhang_bytetrack_2022}. In a dialogue context, robust tracking allows the robot to reliably answer follow-up questions about an entity, sustaining referential continuity.

\section{Semantic Scene Understanding and Scene Graphs}

While bounding boxes localize objects, they lack the relational semantics required for language generation. Scene Graphs~\cite{johnson_image_2015} bridge this gap by providing a structured intermediate representation---a directed graph where nodes represent detected objects and edges represent semantic or spatial relationships (e.g., \verb|<cup> [on] <table_>|).

In addition to pairwise spatial relations, objects possess unary properties such as color or state. In this work, to maintain a uniform graph structure, attributes are elegantly modeled as self-referential edges (unary relations) on the object nodes. This unifies the representation, allowing the dialogue system to query properties and spatial relations using identical graph traversal logic.

Scene Graph Generation (SGG) can be achieved through heuristic geometric rules or learned models. Learned approaches, such as RelTR (Relation Transformer)~\cite{cong_reltr_2023}, treat SGG as a direct set-prediction problem, inferring subject-predicate-object triplets from image features. Reflecting the consensus that hierarchical, persistent representations are central to robust robot understanding~\cite{mascaro_scene_2025, hughes_foundations_2024}, this thesis employs a hybrid neuro-symbolic approach, combining fast geometric rules for stable spatial relations with learned multimodal prompting for richer semantic connections.

\section{Multimodal AI: LLMs, VLMs, and Efficient Architectures}

To generate grounded dialogue, the structured scene representation must be connected to natural language via Large Language Models (LLMs). LLMs are built upon the Transformer architecture, utilizing self-attention mechanisms for autoregressive next-token prediction~\cite{vaswani2023attentionneed}. 

Vision-Language Models (VLMs) extend this capability by conditioning text generation on visual inputs~\cite{liu_visual_2023, dai_instructblip_2023, driess_palm-e_2023}. Typically, a Vision Encoder (such as a ViT) extracts features from an image, which are then projected into the LLM's embedding space. To the LLM, the image functions as a sequence of ``visual words'' appended to the text prompt.

As models scale to achieve better reasoning, computational costs rise dramatically. A dense model activates all its parameters for every token. To mitigate this, modern architectures frequently employ a Mixture of Experts (MoE) routing mechanism~\cite{shazeer2017outrageouslylargeneuralnetworks}. In an MoE model, only a small subset of specialized neural ``experts'' is activated for any given token. Consequently, if a system possesses sufficient RAM or VRAM to load the model weights, it can achieve the reasoning capabilities of a massive parameter model at the inference speed and compute cost of a much smaller model. This architectural efficiency is highly advantageous in robotics, where rapid inference is critical.

Despite their capabilities, VLMs are prone to hallucinating relations or failing to disambiguate multiple instances of the same object. To enforce precise visual grounding, this thesis utilizes Set-of-Mark (SoM) prompting~\cite{yang_set-of-mark_2023}. By rendering explicit visual markers---such as numbered bounding boxes or segmentation masks~\cite{kirillov_segment_2023}---directly onto the image, SoM transforms visual regions into explicit, referable symbols. This forces the VLM to anchor its relational reasoning to specific, tracked detections.

\section{Grounded Dialogue and Retrieval-Augmented Generation (RAG)}

Classical robotic dialogue systems rely on rigid intent-slot architectures~\cite{bocklisch_rasa_2017}. While predictable, these rule-based managers become brittle when users ask unanticipated questions about the environment. Generative AI offers open-ended flexibility, provided the generation can be tethered to the robot's physical reality.

Retrieval-Augmented Generation (RAG)~\cite{lewis_retrieval-augmented_2020} provides a framework for enriching language models with external data. While traditional RAG retrieves static text passages from document databases, this thesis adapts the paradigm for embodied robotics. Here, the ``database'' is the robot's dynamic Scene Memory. The dialogue layer retrieves structured scene graphs, recent visual captions, and object-specific attributes, injecting this live perceptual context directly into the LLM's prompt~\cite{grassi_grounding_2024}.

Crucially, this structured RAG approach shifts the cognitive burden placed on the LLM. Rather than relying on its internal Parametric Memory to hallucinate an answer, the model is tasked with ``reading comprehension'' over the injected prompt context. By providing the exact physical state in the prompt, even relatively small LLMs can perform exceptionally well on grounded Question-Answering tasks. Coupled with the MoE architectures discussed previously, this strategy drastically reduces VRAM requirements and inference latency, enabling highly responsive, factually grounded HRI on computationally constrained platforms.
